% Part: first-order-logic
% Chapter: axiomatic-proofs
% Section: identity

\documentclass[../../../include/open-logic-section]{subfiles}

\begin{document}

\olfileid[pt]{fol}{axd}{ide}

\olsection{\usetoken{P}{derivation} com \usetoken{S}{identity}}

Para incorporar $\eq$ às !!{derivation}s, simplesmente adicionamos novos
esquemas de axiomas. A definição de !!{derivation} e de $\Proves$
permanece a mesma; apenas permitimos também os novos axiomas.

\begin{defn}[Axiomas para a !!{identity}]
\begin{align}
\ollabel{ax:id1} & \eq[t][t], \\
\ollabel{ax:id2} & \eq[t_1][t_2] \lif (!B(t_1) \lif
  !B(t_2)),
\end{align}
para quaisquer termos fechados $t$, $t_1$, $t_2$.
\end{defn}

\begin{prop}\ollabel{prop:sound}
  Os axiomas \olref{ax:id1} e \olref{ax:id2} são válidos.
\end{prop}

\begin{proof}
  Exercício.
\end{proof}

\begin{prob}
Demonstre \olref[fol][axd][ide]{prop:sound}.
\end{prob}

\begin{prop}\ollabel{prop:iden1}
 $\Gamma \Proves \eq[t][t]$, para qualquer termo $t$ e conjunto~$\Gamma$.
\end{prop}

\begin{prop}\ollabel{prop:iden2}
  Se $\Gamma \Proves !A(t_1)$ e $\Gamma \Proves
  \eq[t_1][t_2]$, então $\Gamma \Proves !A(t_2)$.
\end{prop}

\begin{proof}
A !!{formula}
\[
(\eq[t_1][t_2] \lif (!A(t_1) \lif !A(t_2)))
\]
é uma instância de~\olref{ax:id2}. A conclusão decorre de duas aplicações
de~\MP.
\end{proof}

\end{document}
